% from ega4-18-fr.tex
\oldpage[IV-4]{163}

est l'une de ces extensions; donc la fermeture intégrale de $Y$ dans $L$ est somme de préschémas dont l'un est la fermeture intégrale de $Y$ dans $L_1$. Le corollaire résulte donc aussitôt de (18.10.14).

Lorsque par exemple $A$ est l'anneau des entiers d'un corps de nombres algébriques $K$, $K'$ et $L$ des extensions algébriques de $K$, il y a des exemples classiques montrant que la relation $C_1=A[C,A']$ est en défaut lorsque $L$ n'est pas non ramifiée sur $A$ [\ ].

\begin{env}[18.10.15]
\label{IV.18.10.15}
Supposons que $Y=\Spec(A)$ soit affine, $A$ étant intègre et intégralement clos; soient $K$ son corps des fractions, $L$ une extension finie séparable de $K$, et supposons que la fermeture intégrale $C$ de $A$ dans $L$ soit un $A$-module projectif de type fini (ce qui sera par exemple le cas si $A$ est un anneau de Dedekind (Bourbaki, \emph{Alg. comm.}, chap.~VII, \textsection~4, n\textsuperscript{o}~10, prop.~22)). Dire que $L$ est non ramifiée sur $A$ signifie que $C$ est étale (18.10.10) et en vertu de (18.2.7, (ii)), cela équivaut donc à dire que le discriminant $d_{C/A}$ de $\Spec(C)$ sur $\Spec(A)$ est inversible dans $A$. Dans le cas particulier où $C$ est un $A$-module libre et $(x_i)_{1\leq i\leq n}$ une base de $C$ sur $A$, cela signifie que $\det(\operatorname{Tr}_{C/A}(x_ix_j))$ est inversible dans $A$.
\end{env}

\begin{theorem}[18.10.16]
\label{IV.18.10.16}
Soient $Y$ un préschéma intègre, $\eta$ son point générique, $f:X\to Y$ un morphisme séparé et localement de type fini. On suppose que toute composante irréductible de $X$ domine $Y$ et que la fibre générique $X_\eta=f^{-1}(\eta)$ est un préschéma fini sur $K=k(\eta)$, donc $(X_\eta)_{\mathrm{red}}$ est égal à $\Spec(L)$, où $L=\prod_iL_i$ est un produit d'extensions finies $L_i$ de $K$. On pose
\[
  n=[L:K]=\sum_i[L_i:K],\qquad n_s=\sum_i[L_i:K]_s \quad\text{(somme des degrés séparables des $L_i$)}.
\]
Soit $y$ un point géométriquement unibranche de $Y$, et désignons par $n(y)$ la somme des degrés séparables sur $k(y)$ des corps résiduels des points isolés de $f^{-1}(y)$. Alors:
\begin{enumerate}
  \item[(i)] On a $n(y)\leq n_s\leq n$.
  \item[(ii)] Supposons $X$ réduit et le point $y$ normal dans $Y$. Pour que $n(y)=n$, il faut et il suffit qu'il existe un voisinage ouvert $U$ de $y$ dans $Y$ tel que la restriction $f^{-1}(U)\to U$ de $f$ soit un morphisme étale et fini.
\end{enumerate}
\end{theorem}

\emph{A) Réduction au cas où $f$ est quasi-fini.} --- On sait (13.1.4) que l'ensemble $Z$ des points $x\in X$ isolés dans $f^{-1}(f(x))$ est ouvert dans $X$ et par hypothèse $Z$ contient $f^{-1}(\eta)$. Puisque $f^{-1}(\eta)$ est finie, $Z$ est réunion filtrante croissante d'ouverts quasi-compacts $V_\lambda$ contenant $f^{-1}(\eta)$ (donc denses dans $X$ par hypothèse). Si $n_\lambda(y)$ est la somme des degrés séparables sur $k(y)$ des corps résiduels des points de $V_\lambda\cap f^{-1}(y)$, on a $n(y)=\sup_\lambda n_\lambda(y)$; pour prouver (i), il suffira donc de montrer que $n_\lambda(y)\leq n_s$. D'autre part, si $n(y)=n$, il y a un indice $\lambda$ tel que $n_\lambda(y)=n$. Supposons démontré qu'il existe alors un voisinage ouvert $U$ de $y$ dans $Y$ tel que la restriction $V_\lambda\cap f^{-1}(U)\to U$ de $f$ soit un morphisme étale et fini. Puisque $f$ est séparé, l'injection canonique $V_\lambda\cap f^{-1}(U)\to f^{-1}(U)$ est alors aussi un morphisme fini (II, 6.1.5, (v)), donc $V_\lambda\cap f^{-1}(U)$ est fermé dans $f^{-1}(U)$ (II, 6.1.10); mais étant partout dense dans $f^{-1}(U)$, il lui est nécessairement égal.

\emph{B) Réduction au cas où $X$ est réduit, $f$ fini et $Y=\Spec(\sh{O}_{Y,y})$.} --- On notera que l'hypothèse sur $f$ entraîne qu'elle est aussi vérifiée pour $f_{\mathrm{red}}$ et que la conclusion de (i) ne concerne que $f_{\mathrm{red}}$. On peut donc remplacer partout $f$ par $f_{\mathrm{red}}$ et supposer donc $X$ réduit.

